\documentclass[12pt,a4paper]{article}
\usepackage[utf8]{inputenc}
\usepackage[T1]{fontenc}
\usepackage{geometry}
\usepackage{longtable}
\geometry{a4paper, margin=2.5cm}
\title{\textbf{Documento de Requisitos}\\Sistema Reserva de Salas de Estudo}
\author{Ailton Baren Pereira da Silva}
\date{Junho de 2026}
\begin{document}
\maketitle

\section{Objetivo}
Registrar requisitos funcionais, requisitos não funcionais e regras de negócio do sistema Reserva de Salas de Estudo.

\section{Perfis}
\subsection{Aluno}
Consulta salas e horários, cria reservas, visualiza suas reservas e cancela suas próprias reservas.

\subsection{Administrador}
Gerencia usuários, salas, bloqueios, reservas e configurações do sistema.

\section{Requisitos Funcionais}
\begin{longtable}{|p{2cm}|p{12cm}|}
\hline
\textbf{ID} & \textbf{Requisito} \\ \hline
RF-01 & Permitir login de usuários cadastrados. \\ \hline
RF-02 & Permitir logout. \\ \hline
RF-03 & Exibir dashboard conforme o perfil. \\ \hline
RF-04 & Permitir consulta de salas. \\ \hline
RF-05 & Permitir consulta de horários por sala, data e duração. \\ \hline
RF-06 & Permitir criação de reservas com finalidade. \\ \hline
RF-07 & Permitir visualização das próprias reservas. \\ \hline
RF-08 & Permitir cancelamento de reserva com motivo. \\ \hline
RF-09 & Permitir cadastro de usuários pelo administrador. \\ \hline
RF-10 & Permitir edição de usuários pelo administrador. \\ \hline
RF-11 & Permitir ativar e desativar usuários. \\ \hline
RF-12 & Permitir cadastro de salas. \\ \hline
RF-13 & Permitir edição de salas. \\ \hline
RF-14 & Permitir ativar, desativar e bloquear salas. \\ \hline
RF-15 & Permitir visualização administrativa de reservas. \\ \hline
RF-16 & Permitir filtro administrativo de reservas. \\ \hline
RF-17 & Permitir cancelamento administrativo de reservas. \\ \hline
RF-18 & Permitir bloqueio de salas por período. \\ \hline
RF-19 & Permitir cancelamento de bloqueios. \\ \hline
RF-20 & Permitir configuração do limite de reservas ativas por aluno. \\ \hline
RF-21 & Exibir páginas 403 e 404. \\ \hline
\end{longtable}

\section{Requisitos Não Funcionais}
\begin{longtable}{|p{2cm}|p{12cm}|}
\hline
\textbf{ID} & \textbf{Requisito} \\ \hline
RNF-01 & Utilizar Python 3 e Flask. \\ \hline
RNF-02 & Utilizar PostgreSQL como banco da aplicação. \\ \hline
RNF-03 & Utilizar Flask-SQLAlchemy para persistência. \\ \hline
RNF-04 & Utilizar Flask-Login para autenticação. \\ \hline
RNF-05 & Utilizar Jinja2 para templates. \\ \hline
RNF-06 & Utilizar Bootstrap 5 por CDN. \\ \hline
RNF-07 & Ter layout responsivo. \\ \hline
RNF-08 & Possuir testes automatizados com pytest. \\ \hline
RNF-09 & Armazenar senhas com hash seguro. \\ \hline
RNF-10 & Executar localmente sem Docker, React, Node.js, APIs externas ou nuvem. \\ \hline
\end{longtable}

\section{Regras de Negócio}
\begin{longtable}{|p{2cm}|p{12cm}|}
\hline
\textbf{ID} & \textbf{Regra} \\ \hline
RN-01 & A hora final deve ser posterior à hora inicial. \\ \hline
RN-02 & Não aceitar reservas em datas passadas. \\ \hline
RN-03 & A reserva deve ter duração fixa de uma hora. \\ \hline
RN-04 & Uma sala não pode possuir reservas ativas sobrepostas na mesma data. \\ \hline
RN-05 & Somente proprietário ou administrador pode cancelar reserva. \\ \hline
RN-06 & Salas inativas ou bloqueadas não podem ser reservadas. \\ \hline
RN-07 & Um aluno não pode criar reservas sobrepostas. \\ \hline
RN-08 & Somente administradores gerenciam usuários, salas, bloqueios e configurações. \\ \hline
RN-09 & Páginas protegidas exigem autenticação. \\ \hline
RN-10 & Senhas devem ser armazenadas com hash seguro. \\ \hline
RN-11 & Cancelamento de reserva deve registrar motivo. \\ \hline
RN-12 & O administrador pode configurar o limite de reservas ativas. \\ \hline
RN-13 & Bloqueios por período impedem reservas no intervalo. \\ \hline
RN-14 & E-mails de usuários devem ser únicos. \\ \hline
RN-15 & O administrador não pode desativar seu próprio usuário. \\ \hline
\end{longtable}

\section{Dados Iniciais}
O sistema cria administrador, alunos, salas de estudo e configuração inicial de limite de reservas por meio do script de seed.

\end{document}
